\section{Programm: Hallo Welt!}
\label{sec:c:hallo-welt}

In alter Tradition ist das erste Programm, das man beim Erlernen einer neuen Programmiersprache zu sehen bekommt, gewöhnlich das ``Hallo Welt!'' Programm.\footnote{siehe \cite{wikipedia:hallo-welt}} 
Da wir uns zunächst auf die Programmiersprache ``C'' beschränken, ist das:
    
\cppinput{c/01_hallo_welt.c}

\paragraph{Erklärung:}

\begin{description}
    \item[\texttt{/* Das "Hallo Welt"{}-Programm. */}]\hfill\\
        Das ist ein \textbf{Kommentar}.
        Der ``Compiler''\footnote{%
            Der Compiler ist das Programm, das den C-Code in die sogenannte Maschinensprache, also ein Format, das der Prozessor direkt versteht, übersetzt.}
        ignoriert ihn völlig. Er ist nur für uns Menschen da,
        um den Code zu erklären.

        In C gibt es zwei Arten von Kommentaren:
        \begin{itemize}
            \item \texttt{/* ... */} — für längere Kommentare über mehrere Zeilen.
            \item \texttt{// ...} — für kurze Kommentare in einer Zeile.
        \end{itemize}

    \item[\texttt{\#include <stdio.h>}]\hfill\\
        Das ist eine \textbf{Header-Datei} (auf Deutsch auch \emph{Kopfdatei}).

        \texttt{stdio.h} steht für \textit{Standard Input Output}.
        Sie enthält Informationen über Funktionen wie \texttt{printf},
        die wir später benutzen wollen.

        \textbf{Merke:} Wenn wir eine Funktion benutzen wollen, die wir nicht
        selbst geschrieben haben, müssen wir die dazugehörige Header-Datei
        einbinden.

    \item[\texttt{int main() \{ }]\hfill\\
        Das ist die \textbf{Hauptfunktion} (englisch: \textit{main function}).

        Jedes C-Programm startet hier. Der Computer führt den Code in den
        geschweiften Klammern \texttt{\{ ... \}} aus.

        \texttt{int} bedeutet, dass die Funktion eine ganze Zahl (Integer)
        zurückgibt — das ist der Rückgabewert, der weiter unten kommt.

    \item[\texttt{printf("Hallo Welt!\backslash n");}]\hfill\\
        Das ist die \textbf{Ausgabe}.

        \texttt{printf} steht für \textit{print formatted} und schreibt Text
        auf den Bildschirm (genauer: in die Konsole).

        Das \texttt{\backslash n} ist ein \textbf{Zeilenumbruch}.
        Es sorgt dafür, dass der Cursor nach der Ausgabe in die nächste Zeile springt.

        \textbf{Wichtig:} Jede Anweisung in C endet mit einem Semikolon \texttt{;}.

    \item[\texttt{return 0;}]\hfill\\
        Das ist die \textbf{Rückgabe}.

        Sie sagt dem Betriebssystem: \textit{``Alles ist gut gelaufen.''}

        \texttt{0} bedeutet in der Regel \textit{``kein Fehler''}.
        Wenn ein Programm mit \texttt{return 0;} endet, weiß der Computer,
        dass das Programm erfolgreich war.

    \item[\texttt{\}} ]\hfill\\
        Die schließende, geschweifte Klammer beendet die \texttt{main}-Funktion.
        Alles zwischen \texttt{\{} und \texttt{\}} gehört zur Funktion.
\end{description}

\paragraph{Merke:}

\begin{itemize}[noitemsep]
    \item Ein C-Programm beginnt immer in der \texttt{main()}-Funktion.
    \item \texttt{printf()} schreibt Text auf den Bildschirm.
    \item \texttt{\backslash n} macht einen Zeilenumbruch.
    \item \texttt{return 0;} sagt dem Computer: ``Alles OK.''
    \item Jede Anweisung endet mit \texttt{;}.
\end{itemize}

